% Expository adaptation, October 7, 2026. Original mathematics: OpenAI, September 23, 2026.
% Prose revised and explanations expanded; see README.md for the source and scope.
\section{The transition model and resource parameters}\label{sec:model}

We now make the configuration-graph description precise. Two choices in this reduction will be used throughout the proof. Time is part of the vertex, so the transition matrix has no directed cycles. Also, a vertex retains the complete configuration, even after the machine halts. This gives a bounded list of predecessors as well as successors. Merging all accepting configurations into one sink would lose the predecessor bound.

The reduction gives continuation probabilities at every vertex, not only at the designated start. This stronger conclusion is needed because the later construction takes weighted copies of these probabilities.

We use machines with finite control, an endmarked read-only input,
finitely many work tapes, and fresh independent fair bits. Input heads
are confined between the endmarkers. All randomized running-time
bounds hold on every random tape. Space counts traversed writable cells
in bits, including every simultaneously live counter, numerical register
and suspended call; visiting a blank work cell counts toward space.
A transducer has a one-way write-only output tape. Throughout,
$n$ is the input length and logarithms have base two.


A matrix indexed by a finite set $V$ is \emph{strictly forward} if
there is an integer time attached to every vertex and every nonzero
entry goes to a strictly later time. It is \emph{substochastic} if its
entries are nonnegative and each row sums to at most one. An edge is
specified by a vertex and a port from a fixed finite set. A \emph{paired}
edge operation returns the other endpoint together with its inverse
port; applying the inverse operation restores the first endpoint.
Absent ports are explicitly recognized.

\begin{lemma}[Configuration reduction]\label{lem:model-reduction}
Fix a polynomial-time randomized machine using $O(\log(n+2))$ work
space. On an input of length $n$ one can specify, uniformly in
logarithmic space, a polynomial-size vertex set $V_0$, a strictly
forward substochastic matrix $S$, a vector
$e\in\{0,1\}^{V_0}$, and a start vertex $x_0$ with the following
properties.
\begin{enumerate}
\item The incoming and outgoing port degrees of $S$ are bounded by
constants depending only on the machine. Port validity, edge weights,
and pairing are computable in logarithmic space.
\item The vector $p_0=(I-S)^{-1}e$ is well defined and belongs to
$[0,1]^{V_0}$.
\item The coordinate $p_0(x_0)$ is the acceptance probability of
the machine.
\end{enumerate}
In particular, an $\RL$ computation gives the promise
$p_0(x_0)=0$ or $p_0(x_0)\ge1/2$.
\end{lemma}

\begin{proof}
Let $T(n)$ be a fixed polynomial upper bound on the running time.
A vertex records a time in $\{0,\ldots,T(n)\}$, the finite control,
all input and work head positions, and the contents of a sufficiently
large box of $O(\log(n+2))$ work cells. Take elementary machine steps
with fixed dyadic coin probabilities. A transition advances the time
by one. Invalid configurations or moves leaving the box have no
corresponding transition. The box is chosen large enough that this
never affects a computation from the specified start.

We need local access in both directions. Successor access is immediate: each elementary step has constantly many successor choices.
To find predecessors, guess the old finite control, locally overwritten
symbols, head movements, and coin choice, then check that the forward
step gives the specified configuration. There are constantly many such
guesses. Combine transitions with the same ordered pair of endpoints
by adding their probabilities. Order the distinct successors and
predecessors by their configuration encodings; their positions give
paired ports for the resulting matrix entries. These lists have
constant length, so duplicate removal, summed weights, and the inverse
port are computed in logarithmic space by local enumeration. This
base computation makes no recursive graph queries. All records use
$O(\log(n+2))$ bits, and the read-only input can be consulted during
these checks.

The simulated input head is an address stored in a vertex. It must not become an uncounted storage device for the algorithm carrying out the simulation. To ensure this, all input-head coordinates in the vertex universe are restricted to a
fixed polynomial interval containing the positions reachable within
$T(n)$ local steps from their specified initial positions. Configurations
and predecessor candidates outside that interval are invalid. To obtain
a simulated head's scanned symbol, compare its stored coordinate with
the input boundaries, supply any required endmarker or off-data blank
according to the source tape convention, and otherwise scan the real
endmarked input to the designated position using a logarithmic-space
counter. A position forbidden by the source tape convention remains
invalid. The physical scan stays between the real endmarkers, even when
a valid simulated coordinate is outside the data. This routine makes
no recursive graph query.

To collect acceptance at one common time, let an already halted computation keep its complete tape and head data while its time advances. These are distinct vertices in successive layers, not self-loops. Keeping the full configuration prevents many unrelated halted computations from merging at one vertex. Let $e$ be one precisely at accepting
configurations at time $T(n)$, and zero elsewhere. Final-time rows
have no outgoing transitions. Thus $S$ is strictly forward and
substochastic, and it is nilpotent. Its inverse expression is the finite
sum
\[
(I-S)^{-1}e=\sum_{j=0}^{T(n)}S^je.
\]
Starting from any valid configuration, at most one reward can be
collected, at final time. The sum is therefore a probability in
$[0,1]$. At $x_0$ it is exactly the machine's acceptance probability.
\end{proof}

The coordinatewise bound on $p_0$ will be used after weighted projections
to larger vertex sets.  It holds at every configuration, independently
of the acceptance probability at the designated start.  Thus the matrix
construction and the probability approximation below apply to arbitrary
randomized computations satisfying the resource bounds.

\paragraph{The parameters to remember.}
There are three size quantities with different roles. The input has length $n$; a vertex address and the shared randomness will have length $O(B)$; and $N$ is an upper bound on the number of vertices at any one stage. The stage index counts copy steps, whereas a sampling-rate index measures how many vertices a random subset is expected to retain. We specify these quantities below before introducing any of the local thresholds.

\paragraph{Size convention.}
Fix an integer parameter $B$, a power of two, comparable with
$\log_2(n+2)$ and at least a sufficiently large fixed constant.
For example, take the least power of two at least the maximum of that
constant and $\lceil\log_2(n+2)\rceil$.  This choice is computable in
logarithmic space and includes all bounded input lengths in the same
construction. All constants throughout the
construction are independent of the input; they may depend on the fixed
machine and, for fixed inverse-polynomial accuracy, its fixed exponent.
The uniform precision parameter of Theorem~\ref{thm:probability} is
handled separately by padding. The number of stages is $L=O(B)$. At stage $l$, a vertex has
a base identifier followed by $l$ symbols from a fixed radix $D_0$.
The radix will be selected by the hierarchy. Encode each base vertex by
a fixed-width radix-$D_0$ word and prepend the digit $1$; unused base
words are invalid. Appending one copy digit at each stage gives
injective numerical identifiers even across different stage lengths.
Indeed, a base width $w$ places stage-$l$ identifiers in
$[D_0^{w+l},2D_0^{w+l})$, and these intervals are disjoint.
Appending a fixed suffix is an affine map on identifiers.

Every stage has size at most a common bound
\[
N=2^{C_NB}.
\]
We may increase the fixed constant $C_N$ as needed. This is a bound on
one stage, not a claim that there are only $N$ pairs of vertices.
Integrated entry errors below carry a factor $1/N$ while summing over both endpoints. This is a normalization per source vertex, not per ordered pair. Later degree bounds control the additional sum over destinations; inserting a second factor of $1/N$ would change the error statement.

Fix integers $H\ge4$ and $2\le h_0<H$, and put
\[
b_k=2^{-Hk}\qquad(k\ge0).
\]
These are the nominal sampling rates. Stage $l$ will have density
$f=b_l$. The stochastic tables will use rates from this same grid
for their source and terminal retention tests. The error power in
the later analysis is always $s=8$.

\paragraph{Uniform constants and finite ranges.}
All thresholds needed by the algorithm can be chosen rational, with
fixed gaps, and all discrete constants can be chosen integer. We
first prove estimates with constants uniform in the ambient finite
universe. The sampling construction applies to every finite range of
identifiers and rates once its prime field is sufficiently large.
The hierarchy uses $O(B)$ copy digits and rates with $k=O(B)$;
Section~\ref{subsec:final-constants} will verify this rate bound also
for the recursive estimator and give a compatible order of constant
choices.

We can then take the prime field to have size $2^{O(B)}$, with its
exponent chosen after these constants. Its elements and all identifiers
occupy $O(B)$ bits. Increasing that exponent does not change the
probabilistic constants proved in the next section.

